The domain is $-1\le x\le3$. Cauchy-Schwarz gives
\[
\left(x\sqrt{1+x}+\sqrt{3-x}\right)^2
\le(x^2+1)((1+x)+(3-x))=4(x^2+1).
\]
The required positive equality holds exactly when the vectors $(x,1)$ and $(\sqrt{1+x},\sqrt{3-x})$ are proportional with a positive multiplier. In particular $0\le x<3$ and
\[
x\sqrt{3-x}=\sqrt{1+x}.
\]
Squaring under these sign restrictions yields
\[
x^2(3-x)=1+x,\qquad
(x-1)(x^2-2x-1)=0.
\]
Of the roots $1,1-\sqrt2,1+\sqrt2$, only $1$ and $1+\sqrt2$ lie in $[0,3)$. Both satisfy the unsquared proportionality condition. Hence
\[
\boxed{x=1\text{ or }1+\sqrt2}.
\]
